\documentclass{beamer}

\usepackage{comment}
\usepackage{color}
\usepackage{listings}
\usepackage{verbatim}
\usepackage{multicol}
\usepackage{booktabs}
\definecolor{green}{RGB}{0,128,0}

\def\EQ#1\EN{\begin{equation*}#1\end{equation*}}
\def\BA#1\EA{\begin{align*}#1\end{align*}}
\newcommand{\bc}{\begin{center}}
\newcommand{\ec}{\end{center}}
\newcommand{\eq}{\ =\ }

\newcommand\p{\partial}
\newcommand\gD{\mathcal D}
\newcommand{\bnabla}{\boldsymbol{\nabla}}
\newcommand{\bq}{\boldsymbol{q}}
\newcommand{\bOmega}{\boldsymbol{\Omega}}

\newcommand{\caq}{c_\text{aq}}
\newcommand{\gehcomment}[1]{{{\color{orange} #1}}}
\newcommand\add[1]{{{\color{blue} #1}}}
\newcommand\remove[1]{\sout{{\color{red} #1}}}
\newcommand\codecomment[1]{{{\color{green} #1}}}
\newcommand\redcomment[1]{{{\color{red} #1}}}
\newcommand\bluecomment[1]{{{\color{blue} #1}}}
\newcommand\greencomment[1]{{{\color{green} #1}}}
\newcommand\magentacomment[1]{{{\color{magenta} #1}}}

\begin{comment}
\tiny
\scriptsize
\footnotesize
\small
\normalsize
\large
\Large
\LARGE
\huge
\Huge
\end{comment}

\begin{document}
\title{1D Ion Exchange Scenario\ldots\\in a Nutshell}
\author{Glenn Hammond}
\date{\today}

\frame{\titlepage}

%-----------------------------------------------------------------------------
\section{Description of the Voegelin Column}

\frame{\frametitle{Description of the Voegelin Column}

Voegelin et al.\ (2000) measured Ca--Mg--Na ion exchange in soil column C1.
PFLOTRAN applies the Gaines--Thomas model to that experiment.
The simulated domain is a 1D column. The physical column diameter does not enter the model.

\begin{itemize}
  \item Length 48.1 cm, porosity 0.61, pore velocity 5.6 cm/min
  \item Domain: $0.481 \times 1 \times 1$ m, 250 cells along $x$
  \item Darcy velocity: $q = 0.61 \times 5.6$ cm/min $= 5.69333\times 10^{-4}$ m/s
  \item One pore volume: $L\phi/q = 515.36$ s
  \item Inlet changes at 20 and 65 pore volumes
  \item Final time: 80 pore volumes (41228.6 s)
  \item CEC: 71.004 mol/m$^3$ bulk
\end{itemize}

}

%-----------------------------------------------------------------------------
\frame{\frametitle{Cation Exchange Capacity}

The measured CEC is 0.06 mol/kg of soil. PFLOTRAN stores the exchanger capacity per bulk volume,
\[
\omega = \rho_s(1-\phi)\,\mathrm{CEC}.
\]
The grain density implied by the reported mass density per pore volume, $\rho_l = 1.94$ g/cm$^3$, is
\[
\rho_s = \frac{\phi\rho_l}{1-\phi} = 3.0344~\mathrm{g/cm^3}.
\]
With $\phi = 0.61$ this gives $\omega = 71.004$ mol/m$^3$.

}

%-----------------------------------------------------------------------------
\frame{\frametitle{Feed Waters}

Totals are molal. pH is 4.6 in every water.
$\log_{10} P_{\mathrm{CO_2}} = -3.5$.
The inlet tracer steps from $4.65\times 10^{-3}$ to $9.4\times 10^{-3}$ mol/kg.

\footnotesize
\begin{center}
\begin{tabular}{lcccc}
\toprule
 & Initial & Inlet1 & Inlet2 & Inlet3 \\
 & column & 0--20 PV & 20--65 PV & 65--80 PV \\
\midrule
Na$^+$ & $4.65\times 10^{-3}$ & $\sim 0$ & $4.6\times 10^{-3}$ & $4.65\times 10^{-3}$ \\
Ca$^{2+}$ & $5.2\times 10^{-3}$ & $5.3\times 10^{-3}$ & $\sim 0$ & $5.2\times 10^{-3}$ \\
Mg$^{2+}$ & $4.55\times 10^{-3}$ & $\sim 0$ & $2.4\times 10^{-3}$ & $4.55\times 10^{-3}$ \\
\bottomrule
\end{tabular}
\end{center}
\normalsize

Inlet1 is a Ca pulse. Inlet2 replaces Ca with Na and Mg. Inlet3 returns to the initial cations.

}

%-----------------------------------------------------------------------------
\section{Governing Equations}

\frame{\frametitle{Governing Reactive Transport Equations}

\large

\EQ
\frac{\p}{\p t}\left(\phi s\,\Psi_j + \sum_\alpha S_{j\alpha}\right) + \bnabla\cdot\bOmega_j \eq 0
\EN

\EQ
S_{j\alpha} \eq \sum_i \nu_{ji}\, S_{i\alpha}
\EN

\EQ
\bOmega_j \eq \big(\bq - \phi s\,\tau\,\gD\bnabla\big)\Psi_j
\EN

\footnotesize
\begin{align*}
\phi &\eq \text{porosity}; \quad
s \eq \text{liquid saturation}; \quad
\tau \eq \text{tortuosity} \\
\Psi_j &\eq \text{aqueous total component concentration of primary species } j \\
S_{j\alpha} &\eq \text{total sorbed concentration of primary species } j \text{ on site } \alpha \\
S_{i\alpha} &\eq \text{sorbed concentration of surface species } i \text{ on site } \alpha \\
\nu_{ji} &\eq \text{stoichiometry of primary species } j \text{ in surface species } i \\
\bOmega_j &\eq \text{solute flux}; \quad
\bq \eq \text{Darcy velocity} \\
\gD &\eq \text{species-independent diffusion coefficient}
\end{align*}

Flow is not solved. \texttt{SPECIFIED\_VELOCITY} assigns $\bq$.

}

%-----------------------------------------------------------------------------
\frame{\frametitle{Gaines--Thomas Ion Exchange}

\footnotesize
Equivalent fraction of cation $i$ on exchanger $\alpha$, with capacity $\omega$:
\EQ
X_i \eq \frac{z_i S_i}{\omega}, \qquad \sum_i X_i \eq 1.
\EN
With reference cation $j$, the mass-action equation stored by PFLOTRAN is
\EQ
X_i \eq k_i a_i \left(\frac{X_j}{k_j a_j}\right)^{z_i/z_j},
\qquad a_i \eq \gamma_i m_i.
\EN
The \texttt{REFERENCE} cation has $k_j = 1$ and is placed first.
The entered value is the coefficient $k$, not $\log_{10} k$.
For a reference valence $z_j$ and exchange coefficient $K_{ij}$,
\[
k_i = \big(K_{ij}\big)^{1/z_j}.
\]

}

%-----------------------------------------------------------------------------
\frame{\frametitle{Selectivity Coefficients}

Ca$^{2+}$ is the reference cation. The two exchange reactions are
\begin{align*}
2\,\mathrm{Na}^+ + \mathrm{X}_2\mathrm{Ca}
  &\rightleftharpoons \mathrm{Ca}^{2+} + 2\,\mathrm{XNa}, \\
\mathrm{Mg}^{2+} + \mathrm{X}_2\mathrm{Ca}
  &\rightleftharpoons \mathrm{Ca}^{2+} + \mathrm{X}_2\mathrm{Mg}.
\end{align*}

\footnotesize
\begin{center}
\begin{tabular}{lcc}
\toprule
Cation & Entered $k$ & Meaning \\
\midrule
Ca$^{2+}$ & 1 & \texttt{REFERENCE} \\
Na$^+$ & 0.1258925412 & $10^{-0.9}$ \\
Mg$^{2+}$ & 0.6918309709 & $10^{-0.16}$ \\
\bottomrule
\end{tabular}
\end{center}

}

%-----------------------------------------------------------------------------
\section{Input Deck}

\begin{frame}[fragile]\frametitle{SIMULATION}

\begin{itemize}
  \item Subsurface transport only. No flow process model.
  \item \texttt{GIRT}: global implicit reactive transport
\end{itemize}

\begin{semiverbatim}
SIMULATION
  SIMULATION_TYPE SUBSURFACE
  PROCESS_MODELS
    SUBSURFACE_TRANSPORT transport
      MODE GIRT
    /
  /
END

SUBSURFACE
  ...
END_SUBSURFACE
\end{semiverbatim}

\end{frame}

%-----------------------------------------------------------------------------
\begin{frame}[fragile]\frametitle{NUMERICAL\_METHODS}

\begin{itemize}
  \item Direct solver. The column has 250 cells.
  \item Tighten the Newton tolerances to $10^{-8}$.
\end{itemize}

\begin{semiverbatim}
NUMERICAL_METHODS TRANSPORT
  TIMESTEPPER
    TS_ACCELERATION 8
    MAX_STEPS 100000
  /
  NEWTON_SOLVER
    PRECONDITIONER_MATRIX_TYPE AIJ
    RTOL 1.d-8
    ATOL 1.d-8
    STOL 1.d-30
  /
  LINEAR_SOLVER
    SOLVER DIRECT
  /
END
\end{semiverbatim}

\end{frame}

%-----------------------------------------------------------------------------
\begin{frame}[fragile]\frametitle{SPECIFIED\_VELOCITY}

\begin{itemize}
  \item \texttt{SPECIFIED\_VELOCITY} is the Darcy velocity.
  \item $5.69333\times 10^{-4}$ m/s is $0.61 \times 5.6$ cm/min.
  \item The unit string converts the vector to m/s, the internal unit.
\end{itemize}

\begin{semiverbatim}
SPECIFIED_VELOCITY
  UNIFORM? YES
  DATASET 5.69333d-4 0.d0 0.d0 m/s
END
\end{semiverbatim}

\end{frame}

%-----------------------------------------------------------------------------
\begin{frame}[fragile]\frametitle{GRID}

\begin{itemize}
  \item Length 0.481 m, divided into 250 cells ($\Delta x = 1.924$ mm)
  \item Unit cross section. Transport is one-dimensional.
\end{itemize}

\begin{semiverbatim}
GRID
  TYPE STRUCTURED
  NXYZ 250 1 1
  BOUNDS
    0.d0 0.d0 0.d0
    0.481d0 1.d0 1.d0
  /
END
\end{semiverbatim}

\end{frame}

%-----------------------------------------------------------------------------
\begin{frame}[fragile]\frametitle{MATERIAL\_PROPERTY and FLUID\_PROPERTY}

\begin{itemize}
  \item Porosity 0.61 is also stored as \texttt{REFERENCE\_POROSITY}.
  \item Diffusion coefficient $10^{-9}$ m$^2$/s. Tortuosity is 1.
  \item Permeability is unused while the velocity is prescribed.
\end{itemize}

\begin{semiverbatim}
REFERENCE_POROSITY 0.61d0

FLUID_PROPERTY
  DIFFUSION_COEFFICIENT 1.d-9
END

MATERIAL_PROPERTY HD
  ID 1
  POROSITY 0.61
  TORTUOSITY 1.0
  PERMEABILITY
    PERM_ISO 5.43d-13
  /
END
\end{semiverbatim}

\end{frame}

%-----------------------------------------------------------------------------
\begin{frame}[fragile]\frametitle{CHEMISTRY}

\tiny
\begin{semiverbatim}
CHEMISTRY
  PRIMARY_SPECIES
    Na+
    Ca++
    Mg++
    H+
    HCO3-
    Cl-
    Tracer
  /
  SECONDARY_SPECIES
    OH-
    CO3--
    CO2(aq)
    CaOH+
    CaCO3(aq)
    CaHCO3+
    CaCl+
    MgCO3(aq)
    MgHCO3+
    MgCl+
    HCl(aq)
    NaCl(aq)
    NaOH(aq)
  /
  PASSIVE_GAS_SPECIES
    CO2(g)
  /
  MINERALS
    Halite
  /
\end{semiverbatim}

\end{frame}

%-----------------------------------------------------------------------------
\begin{frame}[fragile]\frametitle{CHEMISTRY options}

\begin{itemize}
  \item Hanford database, log formulation, and activity coefficients
  \item \texttt{MOLAL} makes the \texttt{T} constraints molalities
  \item Print free-ion and total molalities for every primary species
  \item Halite is kinetic with a negligible rate. It is not the exchanger.
\end{itemize}

\scriptsize
\begin{semiverbatim}
  MINERAL_KINETICS
    Halite
      RATE_CONSTANT 1.e-30
    /
  /
  DATABASE ../../../database/hanford.dat
  LOG_FORMULATION
  ACTIVITY_COEFFICIENTS
  MOLAL
  OUTPUT
    ALL
    FREE_ION
    TOTAL
  /
END
\end{semiverbatim}

\end{frame}

%-----------------------------------------------------------------------------
\begin{frame}[fragile]\frametitle{ION\_EXCHANGE\_RXN}

\begin{itemize}
  \item CEC is mol/m$^3$ bulk. The \texttt{MINERAL} link is commented out.
  \item Ca$^{2+}$ is \texttt{REFERENCE} and its coefficient is 1.
  \item Na and Mg coefficients are $10^{-0.9}$ and $10^{-0.16}$.
\end{itemize}

\begin{semiverbatim}
  SORPTION
    ION_EXCHANGE_RXN
      CEC 71.004
      CATIONS
        Ca++  1.d0 REFERENCE
        Na+   0.1258925412d0
        Mg++  0.6918309709d0
      /
    /
  /
\end{semiverbatim}

\end{frame}

%-----------------------------------------------------------------------------
\begin{frame}[fragile]\frametitle{REGION and OBSERVATION}

\begin{itemize}
  \item \texttt{west} is the inlet face. \texttt{east} is the outlet face.
  \item The observation at \texttt{east} is the breakthrough curve.
\end{itemize}

\tiny
\begin{semiverbatim}
REGION all
  COORDINATES
    -1.d20 -1.d20 -1.d20
    1.d20 1.d20 1.d20
  /
END
REGION west
  FACE WEST
  COORDINATES
    0. 0. 0.
    0. 1. 1.
  /
END
REGION east
  FACE EAST
  COORDINATES
    0.481 0. 0.
    0.481 1. 1.
  /
END
OBSERVATION
  REGION east
END
\end{semiverbatim}

\end{frame}

%-----------------------------------------------------------------------------
\begin{frame}[fragile,allowframebreaks]\frametitle{CONSTRAINT}

\begin{itemize}
  \item \texttt{T} total molality, \texttt{pH} pH, \texttt{Z} charge balance
  \item \texttt{G} with a non-positive value is $\log_{10}$ of the gas partial pressure in bar
  \item Halite: volume fraction 0.5, area 1 m$^2$/m$^3$
\end{itemize}

\scriptsize
\begin{semiverbatim}
CONSTRAINT Initial
  CONCENTRATIONS
    Na+     4.65d-3  T
    Ca++    5.2d-3   T
    Mg++    4.55e-3  T
    H+      4.6      pH
    HCO3-  -3.5      G   CO2(g)
    Cl-     1.d-3    Z
    Tracer  4.65d-3  T
  /
  MINERALS
    Halite  0.5 1.
  /
END
\end{semiverbatim}

\end{frame}

%-----------------------------------------------------------------------------
\begin{frame}[fragile,allowframebreaks]\frametitle{Inlet constraints}

\scriptsize
\begin{semiverbatim}
CONSTRAINT Inlet1          ! 0 to 20 pore volumes
  CONCENTRATIONS
    Na+     1.d-16  T
    Ca++    5.3d-3  T
    Mg++    1.e-16  T
    H+      4.6     pH
    HCO3-  -3.5     G   CO2(g)
    Cl-     3.d-4   Z
    Tracer  9.4d-3  T
  /
END

CONSTRAINT Inlet2          ! 20 to 65 pore volumes
  CONCENTRATIONS
    Na+     4.6d-3  T
    Ca++    1.d-16  T
    Mg++    2.4e-3  T
    H+      4.6     pH
    HCO3-  -3.5     G   CO2(g)
    Cl-     3.d-4   Z
    Tracer  9.4d-3  T
  /
END

CONSTRAINT Inlet3          ! 65 to 80 pore volumes
  CONCENTRATIONS
    Na+     4.65d-3 T
    Ca++    5.2d-3  T
    Mg++    4.55e-3 T
    H+      4.6     pH
    HCO3-  -3.5     G   CO2(g)
    Cl-     3.d-4   Z
    Tracer  9.4d-3  T
  /
END
\end{semiverbatim}

\end{frame}

%-----------------------------------------------------------------------------
\begin{frame}[fragile]\frametitle{TRANSPORT\_CONDITION}

\begin{itemize}
  \item The list must start at time 0.
  \item 10307.1 s is 20 pore volumes. 33498.2 s is 65 pore volumes.
  \item The outlet reuses the initial water.
\end{itemize}

\scriptsize
\begin{semiverbatim}
TRANSPORT_CONDITION Initial
  TYPE DIRICHLET
  CONSTRAINT_LIST
    0.d0 Initial
  /
END
TRANSPORT_CONDITION west
  TYPE DIRICHLET
  CONSTRAINT_LIST
    0.d0    Inlet1
    10307.1 Inlet2
    33498.2 Inlet3
  /
END
TRANSPORT_CONDITION east
  TYPE DIRICHLET
  CONSTRAINT_LIST
    0.d0 Initial
  /
END
\end{semiverbatim}

\end{frame}

%-----------------------------------------------------------------------------
\begin{frame}[fragile]\frametitle{INITIAL\_CONDITION, BOUNDARY\_CONDITION, STRATA}

\scriptsize
\begin{semiverbatim}
INITIAL_CONDITION Initial
  TRANSPORT_CONDITION Initial
  REGION all
END
BOUNDARY_CONDITION
  TRANSPORT_CONDITION west
  REGION west
END
BOUNDARY_CONDITION
  TRANSPORT_CONDITION east
  REGION east
END
STRATA
  MATERIAL HD
  REGION all
END
\end{semiverbatim}

\end{frame}

%-----------------------------------------------------------------------------
\begin{frame}[fragile]\frametitle{TIME}

\begin{itemize}
  \item Final time 41228.6 s is 80 pore volumes.
  \item The usual maximum step is 20 s.
  \item The step is cut to 1 s across each inlet switch.
\end{itemize}

\scriptsize
\begin{semiverbatim}
TIME
  FINAL_TIME 41228.6 s
  INITIAL_TIMESTEP_SIZE 1. s
  MAXIMUM_TIMESTEP_SIZE 20. s
  MAXIMUM_TIMESTEP_SIZE 1. s at 10200. s
  MAXIMUM_TIMESTEP_SIZE 20. s at 10350 s
  MAXIMUM_TIMESTEP_SIZE 1. s at 33300 s
  MAXIMUM_TIMESTEP_SIZE 20. s at 33600 s
END
\end{semiverbatim}

\end{frame}

%-----------------------------------------------------------------------------
\begin{frame}[fragile]\frametitle{OUTPUT}

\begin{itemize}
  \item Snapshot the column at 20, 65, and 80 pore volumes.
  \item Write the east observation every time step.
  \item Both plotting scripts read \texttt{voegelin-obs-0.pft}.
\end{itemize}

\begin{semiverbatim}
OUTPUT
  TIMES s 10307.1 33498.2 41228.6
  PRINT_COLUMN_IDS
  PERIODIC_OBSERVATION TIMESTEP 1
  FORMAT TECPLOT POINT
END
\end{semiverbatim}

\end{frame}

%-----------------------------------------------------------------------------
\section{Running the Example}

\begin{frame}[fragile]\frametitle{Running PFLOTRAN}

The input deck is \texttt{voegelin.in}.

\begin{semiverbatim}
> cd $PFLOTRAN_DIR
> cd shortcourse/exercises/ion-exchange
> pflotran -input_prefix voegelin
> python3 plot_ionex.py
\end{semiverbatim}

\medskip
\small
Both \texttt{plot\_ionex.py} and \texttt{ionex.gp} read \texttt{exp1.dat}
and \texttt{voegelin-obs-0.pft}.
Python writes \texttt{ionex.png}. To display the curves with gnuplot, run
\texttt{gnuplot -persist ionex.gp}.
Pore volume is time divided by $L\phi/q = 515.36$ s.

\end{frame}

%-----------------------------------------------------------------------------
\begin{frame}\frametitle{Outlet breakthrough}

\begin{figure}
\centering
\includegraphics[width=0.92\linewidth]{../ionex.png}
\end{figure}

\footnotesize
Symbols are the Voegelin et al.\ (2000) measurements.
Lines are PFLOTRAN. Self-sharpening fronts arrive near 10 and 71 pore volumes.
The Mg front from about 30 to 55 pore volumes is self-broadening.

\end{frame}

%-----------------------------------------------------------------------------
\begin{frame}\frametitle{References}

\small
Voegelin, A., V.M. Vulava, F. Kuhnen, and R. Kretzschmar (2000)
Multicomponent transport of major cations predicted from binary adsorption
experiments, \textit{Journal of Contaminant Hydrology}, \textbf{46}, 319--338.

\medskip
DOI: \href{https://doi.org/10.1016/S0169-7722(00)00132-7}{10.1016/S0169-7722(00)00132-7}.

\end{frame}

\end{document}
